

\documentclass{ieeeaccess}
\usepackage{cite}
\usepackage{amsmath,amssymb,amsfonts}
\usepackage{algorithmic}



\usepackage{graphicx}
\usepackage{pgf,tikz}
\usetikzlibrary{babel} %%% <--- Don't forget
\usepackage{textcomp}



%------------------------------------------------------------------------------
% Code for ORCID iD
%------------------------------------------------------------------------------
\usepackage{hyperref,xcolor}
\usepackage{tikz}

%******Error fixing of using '\usepackage{tikz}'******
  \NewSpotColorSpace{PANTONE}
  \AddSpotColor{PANTONE} {PANTONE3015C} {PANTONE\SpotSpace 3015\SpotSpace C} {1 0.3 0 0.2}
  \SetPageColorSpace{PANTONE}
%************************************************


\def\BibTeX{{\rm B\kern-.05em{\sc i\kern-.025em b}\kern-.08em
    T\kern-.1667em\lower.7ex\hbox{E}\kern-.125emX}}




% Make Orcid icon
\definecolor{lime}{HTML}{A6CE39}
\DeclareRobustCommand{\orcidicon}{%
	\begin{tikzpicture}
	\draw[lime, fill=lime] (0,0) 
	circle [radius=0.16] 
	node[white] {{\fontfamily{qag}\selectfont \tiny ID}};
	\draw[white, fill=white] (-0.0625,0.095) 
	circle [radius=0.007];
	\end{tikzpicture}
	\hspace{-2mm}
}
\foreach \x in {A, ..., Z}{%
	\expandafter\xdef\csname orcid\x\endcsname{\noexpand\href{https://orcid.org/\csname orcidauthor\x\endcsname}{\noexpand\orcidicon}}
}

% Define the ORCID iD command for each author separately. Here done for two authors.
\newcommand{\orcidauthorA}{0000-0001-5203-8033}
\newcommand{\orcidauthorB}{0000-0002-8823-2361}
\newcommand{\orcidauthorC}{0000-0003-1705-4411}
%------------------------------------------------------------------------------




\begin{document}
\history{Date of publication xxxx 00, 0000, date of current version xxxx 00, 0000.}
\doi{10.1109/ACCESS.2017.DOI}


\title{Toward Time Series Forecasting in Optimal Sizing of Energy Storage Systems for Hybrid Renewable Energy Systems}

\author{\uppercase{Akanit Kwangkaew}\orcidA{}\authorrefmark{1,2}, %\IEEEmembership{Fellow, IEEE},
\uppercase{Teeradaj Racharak\orcidB{}\authorrefmark{2}, and Chalie Charoenlarpnopparut}.\orcidC{}\authorrefmark{1}
% \IEEEmembership{Member, IEEE}
}
\address[1]{School of Information, Computer, and Communication Technology, Sirindhorn International Institute of Technology, Thammasat University, Pathum Thani, Thailand}
\address[2]{School of Information Science, Japan Advanced Institute of Science and Technology, Ishikawa, Japan}

% This paragraph of the first footnote will contain support 
% information, including sponsor and financial support acknowledgment. 


\tfootnote{``This work was mainly supported by Sirindhorn International Institute of
Technology (SIIT), Thammasat University (TU), and in part by the Japan Advanced Institute of Science and Technology (JAIST), National Science and Technology Development Agency (NSTDA), 
under the JAIST-SIIT dual-degree doctoral program.''}

\markboth
{A. Kwangkaew \headeretal: Toward Time Series Forecasting in Optimal Sizing of ESS for HRES}
{A. Kwangkaew \headeretal:Toward Time Series Forecasting in Optimal Sizing of ESS for HRES}

\corresp{Corresponding author: Akanit Kwangkaew (e-mail: akanitk84@gmail.com).}






\begin{abstract}

%for = because :"He found it increasingly difficult to read, for his eyesight was beginning to fail"
For volatile and uncertain energy sources that generated from renewable energy sources (RES) depend on weather conditions, these sources are dispatched limited. Energy energy storage (ESS) can make the RES more ‘‘dispatch-able and reliable.

% generally causes a dynamic of transmission systems. Hybrid renewable energy systems (HRES), combining RES and traditional energy resources as well as energy storage systems (ESS), need to indicate the optimal battery size in order to yield the highest profit. This research proposes the optimal model determining the size of ESS with optimizing the cost deviation, enabling the state of yielding profit gain for operating systems. Indeed, we focus on determining the optimal size of ESS, which maximizes real-time dispatch profit and reduces the deviation of fluctuating RES subject to forecasting errors. In addition, we detail analysis of the number of window side(s) of forecasting techniques by examining using a simple exponential smoothing (SES) model.


A xxxxx method is proposed to determine optimal size of ESS for volatile energy that generated from RES in electricity markets. The results indicate that the proposed xxxxx and xxxxx method is efficient to find optimal size of ESS. 

\end{abstract}

\begin{keywords}
Simple Exponential Smoothing, SES, Persistence, Predictive modelling
\end{keywords}

\titlepgskip=-15pt

\maketitle

\section{Introduction}
\label{sec:introduction}

% 5 keys of Introduction-------------------

%1.What was I investigating? (problem)
%2.Why was it important? 
%3.What was already know about the subject of my research?
%4.What did I expect to know after doing the research?
%5.How did I  approach the problem?
%-------------------------------------------

%Why ESS is needed.
\PARstart{F}{or} priority areas of clean energy use possesses sufficient quantities of alternative energy to avoid air pollution emission, renewable energy sources (RESs) have increasingly received attention and is by now widely accepted to conventional power systems \cite{b1,b2,b3,b4,b5-wind-data,b6}. Meanwhile, Krabi, which is a southern province of Thailand, is projected to possibly reach around 100\% of electricity consumption using RES in 2026 \cite{b7}. Nevertheless, the characteristics of volatile renewable sources probably lead to an excess of energy production that would be wasted if the demand of consumers and the supply of generated energy does not balance. For instance, a global quantity of curtailed electrical energy of 940.8 billion kWh was calculated in the year 2013 \cite{b9}. Combining renewable and traditional energy resources as well as energy storage systems (ESSs) known as hybrid renewable energy systems (HRESs) is recently appeared to respond with the increasing use of RESs as alternatives to conventional power generation systems.

\par

%1.What was I investigating? (Problem)
% ESS Sizing determination considering the barrier of n long-term investment cost and short-term operational conditions is an important whether any method you use (according to the other researches)
In front of the meter, however, high penetration of volatile renewable sources, which generally depends on weather conditions, probably affects on dispatch-able RES \cite{b10}; thus reducing the deviation is unavoidably needed. To increasing in the power distribution networks, ESSs are suggested to cope with downsides of RESs, which largely depend on the flexibility of resource \cite{b11,b12}. Generally, ESSs have been adopted for mitigating the RES integration problem by reducing the deviation, and for solving other problems such as stabilizing voltage fluctuation (seconds of run-time), smoothing the outputs (minutes of run-time), and providing spinning reserves or ancillary services (hours of run-time) \cite{b14}. However, the barrier for ESSs deployment is optimizing between long-term investment cost and short-term operational conditions \cite{13-main_barrier_ESS}. 

\par
% 2.Why was it important? 
%*** Measuring how far from error distribution is better than measuring the forecasting accuracy in terms of optimal sizing of ESS.



% บอกว่าการตรวจดูforecasting accuracyเพื่อ หา optimal size of ESS อย่างเดียวไม่เพียงพอ ควรที่จะต้องดู character / component ของ time series data set ด้วย


%4.What did I expect to know after doing the research?
%คาดว่าจะนำเสนอ indicator ที่ใช้ในการตวรแต่ละ component ของ time series แยกตามประเภทของ RES โดยเทียบกับ Probability Error Distribution (Mixed distribution function; obtained based on weighted normal and laplace distribution) ที่ได้ตีพิมพ์ไปแล้ว 
%ถ้า component ไหนให้ค่า error distribution ได้ใกล้เคียงกัย Mixed distribution ก็หมายถึงว่ามันให้ benefit ต่อการ optimize size of ESS มากเท่านั้น

%4.What did I expect to know after doing the research?

%*** Time series components (seasonal trend ...) which one important (ranking them)

%*** Measuring each component how far from error distribution is better than measuring the forecasting accuracy in terms of optimal sizing of ESS may contribute the new indication (Just think about how to measure it)


%5.How did I  approach the problem?

\subsection{Subsection}




\appendices

Appendixes, if needed, appear before the acknowledgment.

\section*{Acknowledgment}


% This research is financially supported by Japan Advanced Institute of Science and Technology (JAIST), National Science and Technology Development Agency (NSTDA), Sirindhorn International Institute of Technology (SIIT), Thammasat University (TU) under the JAIST-SIIT dual-degree doctoral program. Also thanks to support from Provincial Electricity Authority (PEA), utility section of Thailand.

% The preferred spelling of the word ``acknowledgment'' in American English is 
% without an ``e'' after the ``g.'' Use the singular heading even if you have 
% many acknowledgments. Avoid expressions such as ``One of us (S.B.A.) would 
% like to thank $\ldots$ .'' Instead, write ``F. A. Author thanks $\ldots$ .'' In most 
% cases, sponsor and financial support acknowledgments are placed in the 
% unnumbered footnote on the first page, not here.




\begin{thebibliography}{00}


% \bibitem{b3} J. U. Duncombe, ``Infrared navigation---Part I: An assessment of feasibility,'' \emph{IEEE Trans. Electron Devices}, vol. ED-11, no. 1, pp. 34--39, Jan. 1959, 10.1109/TED.2016.2628402.

\bibitem{b1} L. Li, Y. Sun, Z. Liu, X. Hou, G. Shi, and M. Su, ``A Decentralized Control with Unique Equilibrium Point for Cascaded-Type Microgrid,” \emph{IEEE Trans. Sustain. Energy}, vol. 10, no. 1, pp. 324-326, 2019, doi: 10.1109/TSTE.2018.2871641.

\bibitem{b2} A. Reinders et al., ``An Exploration of the Three-Layer Model Including Stakeholders, Markets and Technologies for Assessments of Residential Smart Grids,” \emph{Appl. Sci.}, vol. 8, no. 12, p. 23-63, 2018, doi: 10.3390/app8122363.

\bibitem{b3} A. F. Crossland, D. Jones, N. S. Wade, and S. L. Walker, ``Comparison of the location and rating of energy storage for renewables integration in residential low voltage networks with overvoltage constraints,” \emph{Energies}, vol. 11, no. 8, p. 2041, 2018, doi: 10.3390/en11082041.

\bibitem{b4} Y. Sun, X. Hou, J. Yang, H. Han, M. Su and J. M. Guerrero, ``New Perspectives on Droop Control in AC Microgrid," \emph{IEEE Transactions on Industrial Electronics}, vol. 64, no. 7, pp. 5741-5745, July 2017, doi: 10.1109/TIE.2017.2677328.


\bibitem{b5-wind-data} C. Draxl, A. Clifton, B.-M. Hodge, and J. McCaa, The wind integration national dataset (WIND) toolkit, \emph{Applied Energy}, vol. 151, pp. 355366, 2015.

\bibitem{b6} 
N. Uribe-Pérez, L. Hernández, D. de la Vega, and I. Angulo, ``State of the Art and Trends Review of Smart Metering in Electricity Grids,” \emph{Appl. Sci.}, vol. 6, no. 3, pp. 1–24, 2016, doi: 10.3390/app6030068.

% N. Uribe-Pérez, L. Hernández, D. de la Vega, and I. Angulo, “State of the Art and Trends Review of Smart Metering in Electricity Grids,” Appl. Sci., vol. 6, no. 3, pp. 1–24, 2016.


% \item \emph{Basic format for reports:}\\
% \bibitem{b18} R. J. Hijmans and J. van Etten, ``Raster: Geographic analysis and modeling with raster data,'' R Package Version 2.0-12, Jan. 12, 2012. [Online]. Available: \underline {http://CRAN.R-project.org/package=raster} 

% \bibitem{b19} Teralyzer. Lytera UG, Kirchhain, Germany [Online].
% Available:
% \underline{http://www.lytera.de/Terahertz\_THz\_Spectroscopy.php?id=home}, Accessed on: Jun. 5, 2014

\bibitem{b7} ``Krabi Goes Green: Towards A Model Town, With More Than 100\% Renewable Energy,'' Greenpeace Southeast Asia, Thailand, 2018 [Online]. 
Available: \underline{https://www.greenpeace.or.th/report/Krabi-goes-green-EN.pdf}, Accessed on: Oct. 8, 2020


\bibitem{b8} U. Datta, A. Kalam, and J. Shi, “Battery energy storage system for aggregated inertia-droop control and a novel frequency dependent state-of-charge recovery,” \emph{Energies}, vol. 13, no. 8, 2020.

\bibitem{b9} C. Li et al., “Comprehensive review of renewable energy curtailment and avoidance: A specific example in China,” \emph{Renew. Sustain. Energy Rev.}, vol. 41, pp. 1067–1079, 2015.

\bibitem{b10} J. Wu, B. Zhang, H. Li, Z. Li, Y. Chen, and X. Miao, “Statistical distribution for wind power forecast error and its application to determine optimal size of energy storage system,” \emph{Int. J. Electr. Power Energy Syst.}, vol. 55, pp. 100–107, 2014.

\bibitem{b11} 	
N. Jayasekara, M. A. S. Masoum and P. J. Wolfs, ``Optimal Operation of Distributed Energy Storage Systems to Improve Distribution Network Load and Generation Hosting Capability," \emph{IEEE Trans. Sustainable Energy}, vol. 7, no. 1, pp. 250-261, Jan. 2016, doi: 10.1109/TSTE.2015.2487360.
% Jayasekara, N., Member, S., Masoum, M. A. S., Member, S., Wolfs, P. J., and Member, S., “Optimal Operation of Distributed Energy Storage Systems to Improve Distribution Network Load and Generation Hosting Capability.,” pp. 1–12, 2015.


\bibitem{b12} 
R. Fernández-Blanco, Y. Dvorkin, B. Xu, Y. Wang and D. S. Kirschen, ``Optimal Energy Storage Siting and Sizing: A WECC Case Study," \emph{IEEE Trans. Sustainable Energy}, vol. 8, no. 2, pp. 733-743, April 2017, doi: 10.1109/TSTE.2016.2616444.



\bibitem{13-main_barrier_ESS}
J. L. Lorente, X. A. Liu, R. Best, and D. J. Morrow, “Energy storage allocation in power networks - A state-of-the-art review,” \emph{Proc. - 2018 53rd Int. Univ. Power Eng. Conf. UPEC 2018}, 2018.

\bibitem{b14}
H. Bludszuweit, J. A. Dominguez-Navarro and A. Llombart, ``Statistical Analysis of Wind Power Forecast Error," \emph{IEEE Transactions on Power Systems}, vol. 23, no. 3, pp. 983-991, Aug. 2008, doi: 10.1109/TPWRS.2008.922526.

\end{thebibliography}

\begin{IEEEbiography}[{\includegraphics[width=1in,height=1.25in,clip,keepaspectratio]{a1.png}}]{First A. Author} (M'76--SM'81--F'87) and all authors may include 
biographies. Biographies are often not included in conference-related
papers. This author became a Member (M) of IEEE in 1976, a Senior
Member (SM) in 1981, and a Fellow (F) in 1987. The first paragraph may
contain a place and/or date of birth (list place, then date). Next,
the author's educational background is listed. The degrees should be
listed with type of degree in what field, which institution, city,
state, and country, and year the degree was earned. The author's major
field of study should be lower-cased. 

The second paragraph uses the pronoun of the person (he or she) and not the 
author's last name. It lists military and work experience, including summer 
and fellowship jobs. Job titles are capitalized. The current job must have a 
location; previous positions may be listed 
without one. Information concerning previous publications may be included. 
Try not to list more than three books or published articles. The format for 
listing publishers of a book within the biography is: title of book 
(publisher name, year) similar to a reference. Current and previous research 
interests end the paragraph. The third paragraph begins with the author's 
title and last name (e.g., Dr.\ Smith, Prof.\ Jones, Mr.\ Kajor, Ms.\ Hunter). 
List any memberships in professional societies other than the IEEE. Finally, 
list any awards and work for IEEE committees and publications. If a 
photograph is provided, it should be of good quality, and 
professional-looking. Following are two examples of an author's biography.
\end{IEEEbiography}

\begin{IEEEbiography}[{\includegraphics[width=1in,height=1.25in,clip,keepaspectratio]{a2.png}}]{Second B. Author} was born in Greenwich Village, New York, NY, USA in 
1977. He received the B.S. and M.S. degrees in aerospace engineering from 
the University of Virginia, Charlottesville, in 2001 and the Ph.D. degree in 
mechanical engineering from Drexel University, Philadelphia, PA, in 2008.

From 2001 to 2004, he was a Research Assistant with the Princeton Plasma 
Physics Laboratory. Since 2009, he has been an Assistant Professor with the 
Mechanical Engineering Department, Texas A{\&}M University, College Station. 
He is the author of three books, more than 150 articles, and more than 70 
inventions. His research interests include high-pressure and high-density 
nonthermal plasma discharge processes and applications, microscale plasma 
discharges, discharges in liquids, spectroscopic diagnostics, plasma 
propulsion, and innovation plasma applications. He is an Associate Editor of 
the journal \emph{Earth, Moon, Planets}, and holds two patents. 

Dr. Author was a recipient of the International Association of Geomagnetism 
and Aeronomy Young Scientist Award for Excellence in 2008, and the IEEE 
Electromagnetic Compatibility Society Best Symposium Paper Award in 2011. 
\end{IEEEbiography}

\begin{IEEEbiography}[{\includegraphics[width=1in,height=1.25in,clip,keepaspectratio]{a3.png}}]{Third C. Author, Jr.} (M'87) received the B.S. degree in mechanical 
engineering from National Chung Cheng University, Chiayi, Taiwan, in 2004 
and the M.S. degree in mechanical engineering from National Tsing Hua 
University, Hsinchu, Taiwan, in 2006. He is currently pursuing the Ph.D. 
degree in mechanical engineering at Texas A{\&}M University, College 
Station, TX, USA.

From 2008 to 2009, he was a Research Assistant with the Institute of 
Physics, Academia Sinica, Tapei, Taiwan. His research interest includes the 
development of surface processing and biological/medical treatment 
techniques using nonthermal atmospheric pressure plasmas, fundamental study 
of plasma sources, and fabrication of micro- or nanostructured surfaces. 

Mr. Author's awards and honors include the Frew Fellowship (Australian 
Academy of Science), the I. I. Rabi Prize (APS), the European Frequency and 
Time Forum Award, the Carl Zeiss Research Award, the William F. Meggers 
Award and the Adolph Lomb Medal (OSA).
\end{IEEEbiography}

\EOD

\end{document}
